% Part: computability
% Chapter: computability-theory
% Section: introduction

\documentclass[../../../include/open-logic-section]{subfiles}

\begin{document}

\olfileid{cmp}{thy}{int}
\olsection{تعارف}

منطق دی جیہڑی شاخ \emph{قابلِ حسابیت دا نظریہ} کہلاندی اے،
اوہ فنکشناں تے سیٹاں دی قابلِ حسابیت یا نسبتی قابلِ حسابیت
(relative computability) نال متعلق سوالاں نوں گھوکھدی اے۔
کلینی دے اثر دی اک نشانی ایہ اے کہ پہلے ایس موضوع دا ناں
\emph{بازگشت دا نظریہ} سی، تے اج وی دونوں ناں عام ورتے جاندے نیں۔

آؤ فنکشن~$f\colon \Nat \pto \Nat$ نوں \emph{جزوی قابلِ حساب
  (partial computable)} آکھئیے جے اوہنوں حساب دے کسے نمونے وچ
حسابیا جا سکدا ہووے۔ جے $f$ کُل اے، تاں اسی صرف ایہ آکھاں گے
کہ $f$~\emph{قابلِ حساب} اے۔ اک تعلق~$R$ وی قابلِ حساب کہلاندا اے
جے اوہدا خصوصیتی فنکشن~$\Char{R}$ قابلِ حساب ہووے۔ جے $f$
تے~$g$ جزوی فنکشن نیں، تاں اسی $f(x) \fdefined$ ایہ دسّن لئی
لکھاں گے کہ $f$،~$x$ اُتے تعریف شدہ اے، یعنی $x$،~$f$ دے
دائرۂ تعریف وچ اے؛ تے $f(x) \fundefined$ ایس دے اُلٹ معنی لئی،
یعنی ایہ دسّن لئی کہ $f$،~$x$ اُتے تعریف شدہ نہیں۔ اسی
$f(x) \simeq g(x)$ دا مطلب ایہ لاں گے کہ یا تاں $f(x)$ تے
$g(x)$ دونوں غیر تعریف شدہ نیں، یا دونوں تعریف شدہ تے برابر نیں۔

حساب دے کسے خاص نمونے دا حوالہ دِتے بغیر وی قابلِ حسابیت دے
نظریے نوں گھوکھیا جا سکدا اے۔ ایس لئی ایہ وکھایا جاندا اے کہ
اک ہمہ گیر جزوی قابلِ حساب فنکشن~$\fn{Un}(k, x)$ موجود اے۔
ایہ ساڈے لئی جزوی قابلِ حساب فنکشناں دی فہرست بنانا ممکن کردا
اے۔ اسی علامت~$\cfind{k}$ نال $k$-ویں اک آرگومنٹ والے جزوی
قابلِ حساب فنکشن نوں ظاہر کراں گے، جیہدی تعریف
$\cfind{k}(x) \simeq \fn{Un}(k, x)$ اے۔ (کلینی ایس مقصد لئی
$\{ k \}$ ورتدا سی، پر حالیہ وقتاں وچ ایہ علامت اینی عام نہیں
ورتی گئی۔) تھوڑا ہور عام طور اُتے، اسی کسے وی تعداد دے
آرگومنٹاں والے جزوی قابلِ حساب فنکشناں دی فہرست اکّو یکساں
طریقے نال بنا سکدے آں، تے $\cfind{k}[n]$ نال $k$-ویں،
$n$ آرگومنٹاں والے جزوی بازگشتی فنکشن نوں ظاہر کراں گے۔

جے $f(\vec x, y)$ اک کُل یا جزوی فنکشن اے، تاں $\umin{y}{f
(\vec x, y)}$،~$\vec x$ دا اوہ فنکشن اے جیہڑا سب توں
چھوٹا~$y$ واپس دیندا اے جیہدے لئی $f(\vec x, y) = 0$
ہووے، بشرطیکہ $f(\vec x, 0)$، \dots، $f(\vec x, y-1)$
سارے تعریف شدہ ہون؛ جے ایہوجیہا کوئی $y$ نہیں، تاں
$\umin{y}{f (\vec x, y)}$ غیر تعریف شدہ اے۔ جے $R(\vec x, y)$
اک تعلق اے، تاں $\umin{y}{R(\vec x, y)}$ دی تعریف سب توں
چھوٹے~$y$ دے طور اُتے کیتی جاندی اے جیہدے لئی $R(\vec x, y)$
سچ ہووے؛ ہور لفظاں وچ، سب توں چھوٹا~$y$ جیہدے لئی
$1 \tsub \Char{R}(\vec x, y) = 0$ ہووے۔

ایہ وکھاؤن لئی کہ اک فنکشن قابلِ حساب اے، دو طریقیاں نال
اگے ودھیا جا سکدا اے:
\begin{enumerate}
\item سخت ریاضیاتی طریقے نال: اک ٹیورنگ مشین یا جزوی بازگشتی
  فنکشن نوں صاف صاف بیان کرو، تے وکھاؤ کہ اوہ تہاڈے مطلوبہ
  فنکشن دا حساب کردا اے؛
\item غیر رسمی طریقے نال: ایہدا حساب کرن والا اک الگورتھم بیان
  کرو، تے چرچ دے دعوے (Church's thesis) دا سہارا لو۔
\end{enumerate}
دوہاں طریقیاں وچ کوئی تِکھّی حد نہیں؛ الگورتھم دی تفصیلی وضاحت
وچ اینی معلومات ہونی چاہیدی اے کہ ایہ نسبتاً صاف ہووے کہ اصولی
طور اُتے مطلوبہ ٹیورنگ مشین یا جزوی بازگشتی تعریفاں دا سلسلہ
کیویں بنایا جا سکدا اے۔ پوری ریاضیاتی سختی والیاں تعریفاں توں
وڈی وضاحت ملّن دی امید گھٹ اے، تے اسی ایہناں دو طریقیاں وچ
اک مناسب درمیانی راہ لبھن دی کوشش کراں گے؛ مُکدی گل، اسی
ایہوجیہے ثبوت لبھن دی کوشش کراں گے جیہڑے سوجھ نال سمجھ آون
تے ریاضیاتی طور اُتے سخت وی ہون، تے وکھاؤن کہ ٹھیک
ٹھیک تعریفاں حاصل کیتیاں جا سکدیاں نیں۔

\end{document}
